% ============================================================
% Model: 一维方势垒周期势（Kronig-Penney 简化版）
% ============================================================

\section{一维方势垒周期势}
\label{sec:1d-periodic-square-well}

\begin{tipbox}[关键词标签]
周期势 · 傅里叶展开 · 布里渊区边界 · 能隙 · 选择定则 · 近自由电子
\end{tipbox}

% --------------------------------------------------------
% 1. 模型定义框
% --------------------------------------------------------
\begin{modelbox}[模型：一维周期方势垒]
\textbf{物理情境}：一维晶体中电子感受到周期性排列的方势垒（或势阱）的作用。势场在一个周期内分段常数，是最可解析处理的周期势模型之一。

\textbf{势场定义}（周期 $a$）：
\[
V(x)=\begin{cases}
0, & na<x\le(n+\tfrac{1}{2})a,\\
U_0, & (n+\tfrac{1}{2})a<x\le(n+1)a,
\end{cases}\qquad n\in\mathbb{Z}.
\]

\textbf{关键参数}：
\begin{itemize}[nosep]
    \item $U_0$：势垒高度（$U_0>0$ 为势垒，$U_0<0$ 为势阱）
    \item $a$：晶格常数，势场周期
    \item $V_n$：势场的第 $n$ 个傅里叶分量，决定布里渊区边界能隙大小
\end{itemize}

\textbf{边界条件 / 约束}：
\begin{itemize}[nosep]
    \item Bloch 定理：$\psi(x+a)=e^{ika}\psi(x)$
    \item 近自由电子近似：$U_0$ 较小，可用微扰论处理
    \item 能隙出现在 $k=\pm n\pi/a$（布里渊区边界）
\end{itemize}
\end{modelbox}

% --------------------------------------------------------
% 2. 关键公式框
% --------------------------------------------------------
\begin{formulabox}[核心公式]
\begin{enumerate}[label=\textbf{(\arabic*)}]
    \item \textbf{势场傅里叶分量}：
    \[
    V_n=\frac{1}{a}\int_0^a V(x)e^{-iG_n x}\,dx,\qquad G_n=\frac{2\pi n}{a}
    \]

    \item \textbf{方势垒的傅里叶系数}：
    \[
    V_n=\frac{U_0}{\pi n}e^{-i\pi n/2}\sin\frac{\pi n}{2}
    =\begin{cases}
    0, & n\text{ 为偶数},\\
    \displaystyle\frac{U_0}{\pi n}(-1)^{(n-1)/2}, & n\text{ 为奇数}.
    \end{cases}
    \]
    \texteq{选择定则：偶数 $n$ 的傅里叶分量为零，对应能隙消失}

    \item \textbf{布里渊区边界能隙}：
    \[
    E_{g,n}=2|V_n|=\frac{2U_0}{\pi|n|}\quad(n\text{ 为奇数})
    \]
    \texteq{$n=1$ 时 $E_{g,1}=2U_0/\pi$；$n=3$ 时 $E_{g,2}=2U_0/(3\pi)$，依次递减}

    \item \textbf{一般周期势的能隙公式}：
    \[
    E_g(k)=2|V_{G}|,\qquad k=\frac{G}{2}=\frac{n\pi}{a}
    \]
    \texteq{能隙由对应倒格矢的傅里叶分量决定，与势场形状无关}
\end{enumerate}
\end{formulabox}

% --------------------------------------------------------
% 3. 训练点框
% --------------------------------------------------------
\begin{drillbox}[训练点与考法变体]
\trainingpoint{1} \textbf{直接算傅里叶分量}：给定 $V(x)$ 的分段表达式，计算 $V_n$，判断哪些 $n$ 对应非零能隙。

\trainingpoint{2} \textbf{能隙大小排序}：对给定的周期势（正弦、方波、三角波等），比较各布里渊区边界能隙的相对大小。

\trainingpoint{3} \textbf{选择定则分析}：分析势场的对称性（如方势垒的 $V(x+a/2)=U_0-V(x)$ 型对称）如何导致某些 $V_n$ 为零。

\trainingpoint{4} \textbf{Kronig-Penney 精确解}：用波函数在势垒边界处的连接条件（$\psi, \psi'$ 连续），推导色散关系的超越方程。

\trainingpoint{5} \textbf{势垒高度与宽度变化}：固定 $a$，改变 $U_0$ 或势垒宽度 $b$，分析能隙宽度的变化趋势。
\end{drillbox}

% --------------------------------------------------------
% 4. 解题技巧框
% --------------------------------------------------------
\begin{tipbox}[快速识别与解题技巧]
\begin{itemize}
    \item \examnote{识别标志}：题干出现``方势垒''''周期势''''傅里叶分量''''能隙大小''，立即定位本模型。
    \item \examnote{傅里叶积分速算}：
    \[
    \int_{a/2}^{a}e^{-iG_n x}\,dx=\frac{e^{-iG_n a/2}-e^{-iG_n a}}{iG_n}
    =\frac{e^{-i\pi n}-1}{i\cdot 2\pi n/a}=\frac{a}{i2\pi n}\bigl[(-1)^n-1\bigr]
    \]
    当 $n$ 为偶数时分子为零，$n$ 为奇数时值为 $-a/(i\pi n)$。
    \item \examnote{选择定则记忆}：方势垒关于 $x=a/2$ 呈``反中心对称''（$V(x)+V(x+a/2)=U_0$），导致只有奇数次谐波有非零分量。
    \item \examnote{能隙估算}：若 $U_0\ll E_F$，金属性保留；若 $U_0$ 足够大使第一能隙覆盖 $E_F$，则变为绝缘体。
    \item \examnote{量纲检查}：$[V_n]=[E]$，$[E_g]=[E]$，$U_0/\pi$ 具有能量量纲。
\end{itemize}
\end{tipbox}

% --------------------------------------------------------
% 5. 易错点框
% --------------------------------------------------------
\begin{errorbox}{常见失误与陷阱}
\begin{itemize}
    \item \textbf{积分区间}：傅里叶积分是 $\int_0^a$，不是 $\int_{-\infty}^{\infty}$。只在一个周期内积分。
    \item \textbf{$n=0$ 项}：$V_0=U_0/2$ 是势场的平均值，决定能带的整体平移，但不产生能隙。能隙由 $n\neq0$ 的 $|V_n|$ 决定。
    \item \textbf{能隙编号}：$n=1$ 对应第一布里渊区边界 $k=\pi/a$，$n=2$ 对应 $k=2\pi/a$（即 $\Gamma$ 点的第二能带与第一能带之间的能隙？不，是第二与第三之间）。注意：$n$ 为偶数时能隙为零，所以第一个非零能隙是 $n=1$，第二个是 $n=3$。
    \item \textbf{倒格矢与 $n$}：$G_n=2\pi n/a$，布里渊区边界在 $k=G_n/2=n\pi/a$。不要混淆 $n$ 和 $G_n$。
    \item \textbf{KP 模型与近自由电子的区别}：KP 模型可以严格求解（超越方程），近自由电子是微扰近似。当 $U_0$ 较大时，近自由电子的结果可能失效。
\end{itemize}
\end{errorbox}

% --------------------------------------------------------
% 6. 物理图像与关联模型
% --------------------------------------------------------
\begin{insightbox}[物理图像与模型关联]
\textbf{物理图像}：

周期势对电子的作用就像高速公路上的收费站阵列：
\begin{itemize}[nosep]
    \item \textbf{能隙}：收费站的能量壁垒使某些速度（能量）的电子无法传播，形成禁带。
    \item \textbf{选择定则}：如果收费站排列有特殊对称性（如间距的一半处也有结构），某些``谐波''会被抵消，对应能隙消失。
    \item \textbf{势垒高度}：$U_0$ 越大，收费站越难通过，能带越窄（电子越局域）。
\end{itemize}

\textbf{与相似模型的对比}：
\begin{center}
\begin{tabular}{llll}
\toprule
\textbf{对比维度} & \textbf{方势垒周期势} & \textbf{余弦势} & \textbf{Dirac 梳} \\
\midrule
势场形式 & 分段常数 & $V(x)=V_0\cos(2\pi x/a)$ & $V(x)=\sum_n V_0\delta(x-na)$ \\
傅里叶分量 & $V_n\propto 1/n$（奇 $n$） & $V_{\pm1}=V_0/2$，其余为零 & 所有 $V_n=V_0/a$（常数） \\
能隙特征 & 奇 $n$ 有隙，偶 $n$ 无隙 & 只在第一边界有隙 & 所有边界有等大同隙 \\
解析难度 & 中等（分段求解） & 简单（Mathieu 方程） & 简单（$\delta$ 函数） \\
\bottomrule
\end{tabular}
\end{center}

\textbf{推广方向}：
\begin{itemize}[nosep]
    \item 二维/三维：势场的傅里叶分量 $V(\mathbf{G})$ 是矢量，能隙由 $|V_{\mathbf{G}}|$ 决定
    \item 有限深势阱：电子可隧穿通过势垒，能带宽度增大
    \item 准周期势：Fibonacci 序列排列，出现 Cantor 集型的能谱
\end{itemize}
\end{insightbox}

% --------------------------------------------------------
% 7. 推导附录
% --------------------------------------------------------
\begin{derivationbox}[推导细节（自我检验用）]
\textbf{Step 1：傅里叶分量计算}

\begin{align*}
V_n &=\frac{1}{a}\int_0^a V(x)e^{-iG_n x}\,dx
=\frac{U_0}{a}\int_{a/2}^{a}e^{-i\frac{2\pi n}{a}x}\,dx\\
&=\frac{U_0}{a}\cdot\frac{e^{-i2\pi n}-e^{-i\pi n}}{-i\cdot\frac{2\pi n}{a}}
=\frac{U_0}{-i2\pi n}\bigl[1-(-1)^n\bigr].
\end{align*}

当 $n$ 为偶数：$1-(-1)^n=0$，$V_n=0$。

当 $n$ 为奇数：$1-(-1)^n=2$，
\[
V_n=\frac{U_0}{-i\pi n}=\frac{iU_0}{\pi n}=\frac{U_0}{\pi n}e^{i\pi/2}.
\]
取模：$|V_n|=U_0/(\pi|n|)$。

\textbf{Step 2：能隙公式}

在近自由电子近似中，布里渊区边界 $k=n\pi/a$ 处，简并微扰给出能隙
\[
E_g=2|V_n|=\frac{2U_0}{\pi|n|}\quad(n\text{ 为奇数}).
\]

\textbf{Step 3：选择定则的物理根源}

将势场展开为傅里叶级数：
\[
V(x)=\frac{U_0}{2}+\sum_{\substack{n=-\infty\\ n\text{ 奇}}}^{\infty}\frac{U_0}{\pi n}e^{-i\pi n/2}e^{iG_n x}.
\]

方势垒关于 $x=a/4$ 的对称性导致偶数次谐波相消，只有奇数次谐波保留。
\end{derivationbox}

\vspace{1em}
